% GENERATED FILE. Edit canonical lessons, then run npm run generate:books.
% canonical-source-hash: fnv1a64:be423a6b41997e29
% canonical-lessons: RU-C386-doklad, RU-C386-referat, RU-C386-prezentaciya, RU-C386-ponyatie, RU-C386-termin

\chapter{Giving a talk}
\label{ch:ru-giving-a-talk}
\hlchaptermodality{386}

\begin{chapteropening}
\textbf{By the end of this chapter:} \emph{I can give a talk, hand in an essay, prepare a presentation, and explain a concept and a term.}

Complete the last lesson of chapter 386: I can give a talk, hand in an essay, prepare a presentation, and explain a concept and a term.
\end{chapteropening}

\section[\texorpdfstring{\ru{доклад} (doklád)}{доклад (doklád)}]{\ru{доклад} (doklád) — a talk, a paper (given aloud)}

\label{lesson:RU-C386-doklad}



\begin{quote}
Before the new one: say the Russian for a bachelor, then the Russian for a master’s graduate.
\end{quote}

\subsection*{You'll want to know: \ru{доклад}}
\textbf{\ru{доклад}} (\emph{doklád}) — \textquotedblleft{}a talk, a paper (given aloud)\textquotedblright{}.

\begin{quote}
\ru{На} \ru{конференции} \ru{я} \ru{делаю} \ru{доклад}.
\end{quote}

\begin{quote}
\emph{Na konferéntsii ya délayu doklád.}
\end{quote}

\begin{quote}
I am giving a talk at the conference.
\end{quote}

Masculine. A prepared talk, or a report to a superior: \emph{\ru{доклад} \ru{начальнику}}.

\begin{cousinweb}
\textbf{\ru{доклад}} is built on \emph{\ru{докладывать}} \textquotedblleft{}to report\textquotedblright{}, from \emph{\ru{до}-} and \emph{\ru{класть}} \textquotedblleft{}to lay\textquotedblright{}, inherited from Proto-Slavic \emph{*klasti}: to lay a matter before someone.
\end{cousinweb}

\begin{grammarlens}[title={What you've built}]
One more word for telling, explaining and describing at B1.
\end{grammarlens}

\subsection*{Your turn}
\begin{itemize}
\raggedright
  \item \emph{Say it:} \emph{doklád}
  \item \emph{Say it:} \emph{doklád}, once more
  \item \emph{Say it:} a sentence of your own with \emph{doklád}
  \item \emph{Recall:} say the Russian for a bachelor, then the Russian for a master’s graduate, then say \emph{doklád} again
\end{itemize}

\begin{tcolorbox}[breakable,colback=teal!4,colframe=teal!35!black,title={Before you move on}]
What does \ru{доклад} mean? (\textquotedblleft{}A talk, a paper (given aloud)\textquotedblright{}.) Say it once more.
\end{tcolorbox}
\section[\texorpdfstring{\ru{реферат} (referát)}{реферат (referát)}]{\ru{реферат} (referát) — an essay, a summary paper}

\label{lesson:RU-C386-referat}



\begin{quote}
Before the new one: say the Russian for a master’s graduate, then the Russian for a talk.
\end{quote}

\subsection*{You'll want to know: \ru{реферат}}
\textbf{\ru{реферат}} (\emph{referát}) — \textquotedblleft{}an essay, a summary paper\textquotedblright{}.

\begin{quote}
\ru{Студенты} \ru{пишут} \ru{реферат} \ru{по} \ru{литературе}.
\end{quote}

\begin{quote}
\emph{Studénty píshut referát po literatúre.}
\end{quote}

\begin{quote}
The students are writing an essay on literature.
\end{quote}

Masculine. A short paper that sums up what sources say.

\begin{cousinweb}
\textbf{\ru{реферат}} was borrowed from German \emph{Referat} (EtymDB), which is usually traced to Latin \emph{referat} \textquotedblleft{}let him report\textquotedblright{}, a form of \emph{referre} \textquotedblleft{}to bring back, to report\textquotedblright{} (Lewis \& Short).
\end{cousinweb}

\begin{grammarlens}[title={What you've built}]
One more word for telling, explaining and describing at B1.
\end{grammarlens}

\subsection*{Your turn}
\begin{itemize}
\raggedright
  \item \emph{Say it:} \emph{referát}
  \item \emph{Say it:} \emph{referát}, once more
  \item \emph{Say it:} a sentence of your own with \emph{referát}
  \item \emph{Recall:} say the Russian for a master’s graduate, then the Russian for a talk, then say \emph{referát} again
\end{itemize}

\begin{tcolorbox}[breakable,colback=teal!4,colframe=teal!35!black,title={Before you move on}]
What does \ru{реферат} mean? (\textquotedblleft{}An essay, a summary paper\textquotedblright{}.) Say it once more.
\end{tcolorbox}
\section[\texorpdfstring{\ru{презентация} (prezentátsiya)}{презентация (prezentátsiya)}]{\ru{презентация} (prezentátsiya) — a presentation; a launch}

\label{lesson:RU-C386-prezentaciya}



\begin{quote}
Before the new one: say the Russian for a talk, then the Russian for an essay.
\end{quote}

\subsection*{You'll want to know: \ru{презентация}}
\textbf{\ru{презентация}} (\emph{prezentátsiya}) — \textquotedblleft{}a presentation; a launch\textquotedblright{}.

\begin{quote}
\ru{Презентация} \ru{длилась} \ru{двадцать} \ru{минут}.
\end{quote}

\begin{quote}
\emph{Prezentátsiya dlilás dvádtsat minút.}
\end{quote}

\begin{quote}
The presentation lasted twenty minutes.
\end{quote}

Feminine. \emph{\ru{Подготовить} \ru{презентацию}} — to prepare a presentation. Also a launch: \emph{\ru{презентация} \ru{книги}}.

\begin{cousinweb}
\textbf{\ru{презентация}} is usually traced, through western languages, to Late Latin \emph{praesentātiō} \textquotedblleft{}a showing\textquotedblright{}, from \emph{praesentāre} \textquotedblleft{}to show\textquotedblright{}, built on \emph{praesēns} \textquotedblleft{}present, at hand\textquotedblright{} (Lewis \& Short).
\end{cousinweb}

\begin{grammarlens}[title={What you've built}]
One more word for telling, explaining and describing at B1.
\end{grammarlens}

\subsection*{Your turn}
\begin{itemize}
\raggedright
  \item \emph{Say it:} \emph{prezentátsiya}
  \item \emph{Say it:} \emph{prezentátsiya}, once more
  \item \emph{Say it:} a sentence of your own with \emph{prezentátsiya}
  \item \emph{Recall:} say the Russian for a talk, then the Russian for an essay, then say \emph{prezentátsiya} again
\end{itemize}

\begin{tcolorbox}[breakable,colback=teal!4,colframe=teal!35!black,title={Before you move on}]
What does \ru{презентация} mean? (\textquotedblleft{}A presentation; a launch\textquotedblright{}.) Say it once more.
\end{tcolorbox}
\section[\texorpdfstring{\ru{понятие} (ponyátiye)}{понятие (ponyátiye)}]{\ru{понятие} (ponyátiye) — a concept, a notion}

\label{lesson:RU-C386-ponyatie}



\begin{quote}
Before the new one: say the Russian for an essay, then the Russian for a presentation.
\end{quote}

\subsection*{You'll want to know: \ru{понятие}}
\textbf{\ru{понятие}} (\emph{ponyátiye}) — \textquotedblleft{}a concept, a notion\textquotedblright{}.

\begin{quote}
\ru{Я} \ru{не} \ru{имею} \ru{понятия}, \ru{где} \ru{он}.
\end{quote}

\begin{quote}
\emph{Ya ne iméyu ponyátiya, gde on.}
\end{quote}

\begin{quote}
I have no idea where he is.
\end{quote}

Neuter. \emph{\ru{Не} \ru{иметь} \ru{понятия}} — to have no idea.

\begin{cousinweb}
\textbf{\ru{понятие}} is built on \emph{\ru{понять}} \textquotedblleft{}to understand\textquotedblright{}, from \emph{\ru{по}-} and the old verb \textquotedblleft{}to take\textquotedblright{} that also stands in \emph{\ru{снимать}} and \emph{\ru{занимать}}: a concept is what the mind takes hold of.
\end{cousinweb}

\begin{grammarlens}[title={What you've built}]
One more word for telling, explaining and describing at B1.
\end{grammarlens}

\subsection*{Your turn}
\begin{itemize}
\raggedright
  \item \emph{Say it:} \emph{ponyátiye}
  \item \emph{Say it:} \emph{ponyátiye}, once more
  \item \emph{Say it:} a sentence of your own with \emph{ponyátiye}
  \item \emph{Recall:} say the Russian for an essay, then the Russian for a presentation, then say \emph{ponyátiye} again
\end{itemize}

\begin{tcolorbox}[breakable,colback=teal!4,colframe=teal!35!black,title={Before you move on}]
What does \ru{понятие} mean? (\textquotedblleft{}A concept, a notion\textquotedblright{}.) Say it once more.
\end{tcolorbox}
\section[\texorpdfstring{\ru{термин} (términ)}{термин (términ)}]{\ru{термин} (términ) — a term (a technical word)}

\label{lesson:RU-C386-termin}



\begin{quote}
Before the new one: say the Russian for a presentation, then the Russian for a concept.
\end{quote}

\subsection*{You'll want to know: \ru{термин}}
\textbf{\ru{термин}} (\emph{términ}) — \textquotedblleft{}a term (a technical word)\textquotedblright{}.

\begin{quote}
\ru{Этот} \ru{термин} \ru{трудно} \ru{перевести}.
\end{quote}

\begin{quote}
\emph{Étot términ trúdno perevestí.}
\end{quote}

\begin{quote}
This term is hard to translate.
\end{quote}

Masculine. A technical word, not a period of time.

\begin{cousinweb}
\textbf{\ru{термин}} is usually traced, through western languages, to Latin \emph{terminus} \textquotedblleft{}a boundary, a limit\textquotedblright{} (Lewis \& Short): a word with fixed limits to its meaning.
\end{cousinweb}

\begin{grammarlens}[title={What you've built}]
One more word for telling, explaining and describing at B1.
\end{grammarlens}

\subsection*{Your turn}
\begin{itemize}
\raggedright
  \item \emph{Say it:} \emph{términ}
  \item \emph{Say it:} \emph{términ}, once more
  \item \emph{Say it:} a sentence of your own with \emph{términ}
  \item \emph{Recall:} say the Russian for a presentation, then the Russian for a concept, then say \emph{términ} again
\end{itemize}

\begin{tcolorbox}[breakable,colback=teal!4,colframe=teal!35!black,title={Before you move on}]
What does \ru{термин} mean? (\textquotedblleft{}A term (a technical word)\textquotedblright{}.) Say it once more.
\end{tcolorbox}
